\section{Related Work}
\label{sec:related_work}

\noindent \textbf{OSS Supply Chain Security.}
Open-source software (OSS) supply chain security has become a major research focus as attackers increasingly target package registries, dependency ecosystems, and public repositories. Prior works~\cite{gu2023package,ladisa2023sokattack} have systematically studied the attack surface of traditional OSS supply chains, including typosquatting~\cite{ladisa2023sokattack,xie2025squatting}, dependency confusion~\cite{ladisa2023sokattack}, and even package confusion~\cite{neupane2023typosquatting}. In response, researchers have proposed a wide range of detection and mitigation approaches. \emph{Traditional methods} detect suspicious behaviors through signatures~\cite{zheng2024oscar}, semantic rules~\cite{li2023malwukong}, taint-style analysis~\cite{li2023malwukong,duan2021maloss}, and sandbox execution~\cite{zheng2024oscar,duan2021maloss}. \emph{Learning-based methods} use metadata~\cite{sajal2024metadata} and source-code features~\cite{sejfia2022practical,ladisa2023feasibility} to train classifiers for malicious-package detection. More recently, \emph{LLM-based approaches} have been explored for sensitive API identification~\cite{huang2024spiderscan}, data argumentation~\cite{gao2024malguard}, code understanding~\cite{wang2025malpacdetector,yu2024maltracker}, and maliciousness judgment~\cite{zahan2025socketai,di2024posterllm}. In contrast, our work targets a different supply-chain unit, agent skills, whose malicious behavior is often distributed across heterogeneous artifacts rather than only contained in code scripts.

\noindent \textbf{Agentic Supply Chain Security.}
With the rise of LLM agents, recent work has started to examine security risks in the agentic supply chain~\cite{hu2025llmsc,wang2025llmsc,huang2025llmsc}, including agent skills~\cite{liu2026maliciousagentskillswild,liu2026vulskill,guo2026skillprobesecurityauditingemerging} and MCP servers~\cite{hou2026mcp,wang2025mcpguardautomaticallydetecting}. Public reports and benchmark efforts have shown that real-world skill ecosystems already contain attacks such as credential theft~\cite{koi2026clawhavoc}, wallet theft~\cite{koi2026clawhavoc,opensourcemalware-clawdbot-crypto}, reverse shells~\cite{snyk2026skills}, prompt injection~\cite{snyk2026skills}, and agent hijacking~\cite{liu2026maliciousagentskillswild}. Preliminary defenses include rule-based scanners~\cite{huifer-skill-security-scan,goplus-agentguard}, LLM-based analysis~\cite{cisco-skill-scanner}, and dynamic analysis~\cite{liu2026maliciousagentskillswild} for skills. However, most existing approaches analyze individual artifacts in isolation and provide limited support for cross-artifact reasoning. \toolname addresses this gap by focusing on neuro-symbolic reasoning on heterogeneous artifacts.